\par\SourceLeaf{134}
Precisely this, however, gives a very troubling character to the lack of clarity\footnote{Anyone who takes exception to this description of an arrangement which has been in force for half a century is asked to suspend judgment until he has read the whole first division.} of the regulations in the provision intended to govern the amount and progression of the pension rate. The actual words of the rescript of the University Board of Curators dated 20 November 1794, which sanctioned this rule, are as follows:
\begin{quote}
— — — — Second, we approve that whenever the capital has increased by 5000 Thl., and the other revenues of the widows' fund have suffered no decrease, and the number of pensions has not risen above 15, each pension shall be increased by 10 Thl.
\end{quote}
\noindent\hspace*{\parindent}This provision is defective in that it does not make entirely unambiguous \textit{what} it is about a pension increase that is to depend on the preceding conditions: its continuation, or merely its beginning; in other words (considering only the last condition),
\begin{quote}
whether an increase which, in accordance with the statute, began at a time when the number of widows was at most 15 is to cease again, or at least be suitably modified, as soon as the number of pensions later exceeds that standard number,
\end{quote}
or else
\begin{quote}
whether an increase, once begun, is nevertheless to continue unaltered, independently of the standard number being exceeded later.
\end{quote}
\noindent\hspace*{\parindent}I am all the less inclined to dispute that the above formulation of the rule leads more naturally to the second interpretation than to the first, if one considers only the wording without any regard for the consequences that emerge from considering its substance, since the linguistic standpoint is quite subordinate in my present investigation, directed chiefly towards the substance of the arrangement itself. I cannot, however, refrain from setting down here for comparison the form in which \textsc{Brandes} cites the decree in his well-known work on the University of Göttingen. Since \textsc{Brandes} served in 1794 as the ministry's official responsible for university affairs, his understanding is certainly relevant to the matter, even though, from a legal standpoint (likewise outside my investigation), it cannot supersede the wording once used in the official rescript. The passage cited, p. 254, reads:
\begin{quote}
It was further agreed at that time that whenever the fund's capital had grown by 5000 Thl., each pension should be increased by 10 Thl., \textit{as long as} the number of pensioners does not exceed 15, which has never yet been the case.
\end{quote}
\noindent\hspace*{\parindent}Accordingly, \textsc{Brandes} seems to have understood the proviso concerning the 15 pensions rather in the sense of the \textit{first} interpretation, but to have regarded it as wholly irrelevant in practice because of the circumstance he mentions last. This may perhaps also explain why the wording of the rescript was not framed with the utmost precision.

\textsc{Brandes} would, however, probably have judged quite differently had he first examined the facts carefully. In saying `which has never yet been the case', \textsc{Brandes} cannot have meant the seven-year period from the introduction of the provision to the writing of his book, since no weight whatever could be attached to experience from so short a period; he must have meant the whole time since the fund was established. But then his assertion is factually incorrect. The number of pensions paid by the fund had indeed once risen above 15, namely during the first six months of 1778, when there were 16 pensions; indeed, the same would already have happened
\SourceLeaf{135}
two years earlier if the extension of orphans' pensions introduced only in 1794, prolonging them to the age of 20, had applied at that time. This circumstance appears very significant when one considers (in accordance with the remark already made above) that 20—30 years before that period, namely in 1747—1757, the average number of participants was only 21—22, whereas in 1794—1802 it was 35—36.

Subsequently, the number of pensions did not again rise above 15 until Easter 1837, and it has remained above that number ever since. The fund has continued paying the full pension amount, including all the successive increases since 1794, and has thus accepted the second interpretation of the doubtful point at least \textit{implicitly} — for I do not know whether the Church Deputation had previously deliberated on the matter.

But if, setting aside linguistic and formal legal considerations, we wish to decide between the two interpretations solely on their intrinsic grounds, the first question that immediately presents itself is this: if there is really no cause for concern in allowing an increased pension to \textit{continue} undiminished even for 16 recipients, why should it be forbidden to let it \textit{begin} even with 16 recipients? The danger, if there is any, is exactly as great in the one case as in the other. This argument becomes still more striking when we proceed to larger numbers. For allowing an increased pension to continue with 17 recipients (and how much more with 18, 19, and so forth) is quite evidently more dangerous to the fund than allowing the increase to begin with 16; yet, if the second interpretation is adopted, the former is permitted and the latter forbidden.

It is easy to see what conclusions follow from a strictly consistent application of the second interpretation.

In the printed regulations of 1838, §. 10 introduces the arrangement under discussion with the words: The widows' pension \textit{is} \textit{determined} partly by the amount of the fund's assets, partly by the number of widows. According to the introduction prefixed to the first publication of the regulations (1833), they are intended to establish the participants' obligations and rights, and a new member joins by signing a copy presented to him so that he may acquaint himself with those duties and rights. One must therefore assume that the regulations contain them \textit{completely}, and that the fund cannot assert against its members or their survivors any rights not contained in these regulations.

But they contain not a single word reserving the right to reduce the pension rate again if necessary, unless one accepts the first interpretation of the rule determining that rate according to the number of widows. Otherwise, strict consistency requires one to concede that the fund is obliged to continue paying undiminished all pension increases that came into effect while the number of pensions was under 16, even if the number of widows (and orphans) should later (for example, as a result of the overall extension of participation) \textit{grow to any size whatsoever.}

But this is downright absurd unless one assumes that the government is required to provide a guarantee, which is beyond my competence to judge.

